\documentclass[11pt,letterpaper]{article}
\PassOptionsToPackage{hyphens}{url}
\usepackage[margin=.83in,top=.8in,bottom=.8in,headheight=15pt]{geometry}
\usepackage[T1]{fontenc}
\usepackage{lmodern,amsmath,amssymb,mathtools,bm,booktabs,array,graphicx,tikz}
\usetikzlibrary{arrows.meta,positioning,calc,decorations.pathreplacing,patterns}
\definecolor{ink}{HTML}{17344D}\definecolor{accent}{HTML}{087E8B}\definecolor{pale}{HTML}{EDF5F5}\definecolor{gold}{HTML}{AD6B19}
\providecommand{\dd}{\,\mathrm{d}}\providecommand{\R}{\mathbb{R}}
\providecommand{\E}{\mathbb{E}}\providecommand{\Var}{\operatorname{Var}}
\providecommand{\Cov}{\operatorname{Cov}}\providecommand{\diag}{\operatorname{diag}}
\providecommand{\norm}[1]{\left\lVert#1\right\rVert}
\providecommand{\abs}[1]{\left|#1\right|}
\usepackage{microtype,hyperref,fancyhdr,enumitem,amsthm,sectsty,needspace}
\allsectionsfont{\raggedright}
\hypersetup{colorlinks=true,linkcolor=ink,urlcolor=accent,pdfauthor={Zan Ahmad},pdftitle={Membrane current balance and voltage-dependent gates / Review}}
\setlength{\parindent}{0pt}\setlength{\parskip}{6pt}\linespread{1.07}
\setlist{itemsep=3pt,topsep=4pt,leftmargin=*}
\emergencystretch=2em\clubpenalty=10000\widowpenalty=10000
\pagestyle{fancy}\fancyhf{}\fancyhead[L]{\small Computational Cardiology}\fancyhead[R]{\small Lecture 9}
\fancyfoot[L]{\footnotesize Zan Ahmad}\fancyfoot[R]{\footnotesize\thepage}\renewcommand{\headrulewidth}{.3pt}
\newtheorem{proposition}{Proposition}\newtheorem{theorem}{Theorem}\newtheorem{lemma}{Lemma}
\theoremstyle{definition}\newtheorem{definition}{Definition}\newtheorem{example}{Example}
\numberwithin{equation}{section}
\begin{document}
\section*{Membrane current balance and voltage-dependent gates --- review sheet}
Use this after attempting the worked examples and practice. Explain each claim in words before calculating.

\section*{Key ideas}
\begin{itemize}
\item Current signs and density units determine the voltage balance.
\item The gate equilibrium depends on voltage; rest is a coupled current-balance root.
\item Removable rate singularities require finite continuous evaluation.
\item Exact frozen substeps preserve admissibility but do not remove coupling error.
\item Neuronal and cardiac cell models answer different biological questions.
\end{itemize}
\section*{Common errors and repairs}
\textbf{Clipping an unstable gate to zero or one.} Use an admissible update and diagnose the step or equation; clipping changes the dynamics.

\textbf{Calling any spike-shaped trace a cardiac action potential.} Name the model, parameter convention and intended observable.

\textbf{Changing pulse duration inadvertently during refinement.} Treat pulse edges as timeline events and audit the integral.

\section*{Retrieve without your notes}
\begin{itemize}
\item What are the units of alpha and beta?
\item What does zero net current establish at gate equilibrium?
\item Why compare with an independent ODE solver?
\end{itemize}
\section*{Readiness check}
\begin{itemize}
\item Derive the voltage-current balance with consistent units and outward-positive ionic currents.
\item Implement a frozen-rate gate update and explain its interval-preserving property.
\item Evaluate removable rate singularities and compute a coupled resting state.
\item Compare stimulus responses and test time accuracy without equating bounded gates with an accurate trajectory.
\item Distinguish a neuronal HH cell from a reduced cardiac cell model.
\end{itemize}
\end{document}
